% Target: rmp-iii-a-torus-three
% Formula: Moving O_g through the PEPS past O_h inserts the rails g^{-1} and g.
% Ink: Kernel flat frame; the stock MPO skin is the signed box declaration (#5362).
% Boundary: The PEPS bonds and both MPO strings remain open at the window.
% Source: arXiv:2011.12127, TN-Review-main.tex lines 1602-1606; author source
%   ImagesReview Section III/ImagesReview Section III.tex lines 572-603.
% Capabilities: kernel, peps-window, crossing-tensor, string-slide
\[ \tenkzkernel \begin{tenkzeq}
  \begin{tenkz}[rows={wire,wire,wire,wire,wire,wire}, cols=5, bonds=none]
    % The black bonds form the open 5x5 PEPS window.
    \tnwire[name=h1]{(2,1)}{(2,5)} \tnwire[name=h2]{(3,1)}{(3,5)}
    \tnwire[name=h3]{(4,1)}{(4,5)} \tnwire[name=h4]{(5,1)}{(5,5)}
    \tnwire[name=h5]{(6,1)}{(6,5)} \tnwire[name=v1]{(1,1)}{(6,1)}
    \tnwire[name=v2]{(1,2)}{(6,2)} \tnwire[name=v3]{(1,3)}{(6,3)}
    \tnwire[name=v4]{(1,4)}{(6,4)} \tnwire[name=v5]{(1,5)}{(6,5)}
    % Initially O_g lies above the PEPS, while O_h crosses the window.
    \tnwire[kind=string, species=operator, name=g,
      cross={over at crossing of g and v1, over at crossing of g and v2,
             over at crossing of g and v3, over at crossing of g and v4,
             over at crossing of g and v5, over at crossing of g and h}]
      {0.2 w of (1,1)}{0.2 e of (1,5)} \tn[species=operator, at=crossing of g and v1]{}
    \tn[species=operator, at=crossing of g and v2]{}
    \tn[species=operator, at=crossing of g and v3]{}
    \tn[species=operator, at=crossing of g and v4]{}
    \tn[species=operator, at=crossing of g and v5, label pos=0]{O_g}
    \tnwire[kind=string, species=passive, name=h,
      cross={over at crossing of h and h1, over at crossing of h and h2,
             over at crossing of h and h3, over at crossing of h and h4,
             over at crossing of h and h5}]
      {midway (1,3) and (1,4)}{midway (6,3) and (6,4)}
    \tn[species=passive, at=crossing of h and h1]{}
    \tn[species=passive, at=crossing of h and h2]{}
    \tn[species=passive, at=crossing of h and h3]{}
    \tn[species=passive, at=crossing of h and h4]{}
    \tn[species=passive, at=crossing of h and h5, label pos=270]{O_h}
    % Downward pull of the MPO chain; the middle leg carries the arrowhead.
    \tn[at=0.6 w of (2,1), skin=none, name=jp1]{}
    \tn[at=0.6 w of (3,1), skin=none, name=jp2]{}
    \tnwire[kind=string, route=arc, name=p1]{0.4 w of (1,1)}{jp1}
    \tnwire[kind=string, route=arc, dir=to, name=p2,
      cross={over at crossing of p1 and p2}]{jp1}{jp2}
    \tnwire[kind=string, route=arc, name=p3,
      cross={over at crossing of p2 and p3}]{jp2}{0.4 w of (4,1)}
  \end{tenkz} = \begin{tenkz}[rows={wire,wire,wire,wire,wire,wire}, cols=5, bonds=none]
    \tnwire[name=h1]{(2,1)}{(2,5)} \tnwire[name=h2]{(3,1)}{(3,5)}
    \tnwire[name=h3]{(4,1)}{(4,5)} \tnwire[name=h4]{(5,1)}{(5,5)}
    \tnwire[name=h5]{(6,1)}{(6,5)} \tnwire[name=v1]{(1,1)}{(6,1)}
    \tnwire[name=v2]{(1,2)}{(6,2)} \tnwire[name=v3]{(1,3)}{(6,3)}
    \tnwire[name=v4]{(1,4)}{(6,4)} \tnwire[name=v5]{(1,5)}{(6,5)}
    % After the pull, O_g runs below the crossing tensor.
    \tnwire[kind=string, species=operator, name=g,
      cross={over at crossing of g and v2, over at crossing of g and v3,
             over at crossing of g and v4, over at crossing of g and h}]
      {midway (4,1) and (5,1)}{midway (4,5) and (5,5)}
    \tn[species=operator, at=midway (4,1) and (5,1)]{}
    \tn[species=operator, at=crossing of g and v2]{}
    \tn[species=operator, at=crossing of g and v3]{}
    \tn[species=operator, at=crossing of g and v4]{}
    \tn[species=operator, at=midway (4,5) and (5,5), label pos=0]{O_g}
    % O_h remains vertical.  The crossing is the enlarged MPO tensor.
    \tnwire[kind=string, species=passive, name=h,
      cross={over at crossing of h and h1, over at crossing of h and h2,
             over at crossing of h and h3, over at crossing of h and h4,
             over at crossing of h and h5}]
      {midway (1,3) and (1,4)}{midway (6,3) and (6,4)}
    \tn[species=passive, at=crossing of h and h1]{}
    \tn[species=passive, at=crossing of h and h2]{}
    \tn[species=passive, at=crossing of h and h3]{}
    \tn[skin=mpo, role=marked, at=crossing of g and h]{}
    \tn[species=passive, at=crossing of h and h4]{}
    \tn[species=passive, at=crossing of h and h5, label pos=270]{O_h}
    % The two rails conjugate h by g; each marker anchors its rail wires.
    \tn[species=operator, at=0.2 e of (2,3), label pos=90, name=giA]{g^{-1}}
    \tn[species=operator, at=0.2 e of (3,3), name=giB]{}
    \tn[species=operator, name=l3, at=0.2 e of (4,3)]{}
    \tn[species=operator, at=0.2 w of (2,4), label pos=90, name=grA]{g}
    \tn[species=operator, at=0.2 w of (3,4), name=grB]{}
    \tn[species=operator, name=r3, at=0.2 w of (4,4)]{}
    \tnwire[kind=string, species=operator, name=gi1,
      cross={over at crossing of gi1 and h1}]{0.2 e of (1,3)}{giA}
    \tnwire[kind=string, species=operator, name=gi2,
      cross={over at crossing of gi2 and h1, over at crossing of gi2 and h2}]{giA}{giB}
    \tnwire[kind=string, species=operator, name=gi3,
      cross={over at crossing of gi3 and h2, over at crossing of gi3 and h3}]{giB}{l3}
    \tnwire[kind=string, species=operator, name=gr1,
      cross={over at crossing of gr1 and h1}]{0.2 w of (1,4)}{grA}
    \tnwire[kind=string, species=operator, name=gr2,
      cross={over at crossing of gr2 and h1, over at crossing of gr2 and h2}]{grA}{grB}
    \tnwire[kind=string, species=operator, name=gr3,
      cross={over at crossing of gr3 and h2, over at crossing of gr3 and h3}]{grB}{r3}
    % The cup's own via point sits exactly on h's path (deliberate), so it
    % becomes a junction atom; its far end already is gr's endpoint r3.
    \tn[at=on h 0.65, skin=none, name=jc1]{}
    \tnwire[kind=string, route=arc, species=operator, name=c1,
      cross={over at crossing of c1 and h,
             over at crossing of c1 and gi3}]{0.2 e of (4,3)}{jc1}
    \tnwire[kind=string, route=arc, species=operator, name=c2,
      cross={over at crossing of c1 and c2, over at crossing of c2 and gr3}]{jc1}{r3}
  \end{tenkz} \end{tenkzeq}. \]
